\section{2003: Definitive Constitutions, and What Article 8 Says}

\subsection{The approval}

On 29 June 2003 the Pontifical Commission \work{Ecclesia Dei}
definitively approved and confirmed the constitutions of the Priestly
Fraternity of Saint Peter. The approving rescript, protocol 108/2003, is
printed in the Fraternity's own \work{Vademecum}; the canonical
literature quotes its Latin. It records that the Commission acts at the
request of the superior general, Father Arnaud Devillers, in virtue of
faculties granted it by John Paul II; that the text approved is the text
``amended by the General Chapter of this Institute held in the year 2000
and proposed for the definitive approval of the Holy See''; that the three
years of prolonged experimentation have elapsed; and that in printing the
constitutions mention is to be made of the approval granted by the
Apostolic See.\endnote{Pontificia Commissio \work{Ecclesia Dei},
\latin{Rescriptum approbationis Constitutionum ``Fraternitatis
Sacerdotalis Sancti Petri''}, 29 June 2003, Prot.\ N.~108/2003, printed in
the Fraternity's \work{Vademecum} (Fribourg, 2021), 33, and quoted in
Latin by Pietras, art.\ cit., 176--177 n.~20. Bounded negative: the
\work{Vademecum} was not reached for this account, the rescript was not
located in \work{Acta Apostolicae Sedis} or on vatican.va, and the
complete text of the constitutions has never been published. The same
note records that the constitutions were modified by the general chapters
of 2012 and 2018 with the approval of the Apostolic See under Prot.\
N.~153/2009.}

Two things in that record deserve notice. The text definitively approved
in 2003 is the text the chapter of 2000 amended---the chapter that had
just been deprived of its election and had lodged a recourse. And the
intervening sequence is not public: the constitutions were approved
\latin{ad experimentum} for three years in October 1988, reviewed in
detail by the Holy See in 1999, revised by the chapter of 2000, and
definitively approved in 2003, but the instruments that prolonged the
experimental approval across the intervening twelve years were not located
for this account and are not described here.\endnote{Bounded negative. The
Fraternity's own presentation of the 2003 excerpt states that ``until that
date the constitutions had been approved on an experimental basis. They
were reviewed in detail by the Holy See in 1999 and by the Fraternity's
General Chapter (the supreme governing body of the institute) in 2000.''
No instrument of prolongation between 1991 and 2003 was located.}

\subsection{The published excerpt}

What the Fraternity publishes is not the constitutions but an excerpt:
Part I, ``Nature, spirit and aim of the Fraternity,'' articles 1 to 15,
in an English rendering of a French original. Nothing beyond article 15 is
public, which means that the articles on government, on incorporation, on
formation, on the general chapter and on the relations with diocesan
bishops---the articles that would answer half the questions this account
has to leave open---cannot be examined here.\endnote{``Excerpt of the
Constitutions of the Priestly Fraternity of Saint Peter,'' 29 June 2003,
English text published by the Fraternity (fssp.org, ``Documents /
Foundation''), marked ``[Original: French],'' read 2026-07-25. Bounded
negative: articles 16 and following are not published, and the numbering
of article 19, cited in the recourse of July 2000 as providing that the
number of terms of a superior general is not limited, cannot be checked
against the definitive text. The excerpt itself is the definitive text as
approved in 2003, not the 1988 text.}

Article 1 states the canonical figure in the words of canon 731 §\,1: the
Fraternity is ``a clerical society of Apostolic Life of Pontifical Right
whose members, according to the terms of the law, pursue the proper
apostolic aim of the society, and, by leading a communal life according
to a specific form, strives towards the perfection of charity by the
observance of the Constitution.''

Article 6 states the ecclesial bond and cites two councils in one
sentence: founded ``in the spirit of the apostolic letter Motu proprio
Ecclesia Dei Adflicta,'' the Fraternity ``professes its fidelity to the
Roman Pontiff, who, according to the words of the First Vatican Council
(\work{Pastor æternus}), is the `successor of Blessed Peter, Prince of
the Apostles, Vicar of Christ, head of the whole Church, the Father and
Doctor of all Christians' (see Vatican Council II, \work{Lumen Gentium},
n.~22). `Each one of its members is held to obey the Supreme Pontiff as
his highest superior' (Can.\ 590, §\,2).''

Article 7 gives the object: ``the sanctification of priests through the
exercise of the priesthood, and in particular, to turn the life of the
priest toward that which is essentially his \latin{raison d'être}, the
Holy Sacrifice of the Mass.''

And article 8 gives the particular aim, in the sentence that has carried
the whole weight of every liturgical controversy since:

\begin{quote}\small
The particular aim of the Fraternity of Saint Peter is to achieve this
objective through the faithful observance of the ``liturgical and
disciplinary traditions'' according to the dispositions of the Motu
proprio Ecclesia Dei of July 2, 1988, which is at the origin of its
foundation.
\end{quote}

\subsection{Reading article 8 exactly}

Three features of that sentence repay attention.

\emph{It says ``faithful observance,'' not ``exclusive use.''} The word
``exclusive'' does not appear. It appears in the 1999 circular of the
superior general, who speaks of a preferential option to maintain the
exclusive use of the older liturgy in the Fraternity's ministry; and it
appears in the canonical literature, where the 2000 recourse's ``specific
element'' is glossed as ``the exclusive celebration of the liturgy
according to the books of the \latin{editio typica} of 1962.'' It does not
appear in the recourse's own words, and it does not appear in the article
that is supposed to be its source.\endnote{The circular of 30 August 1999
is cited above from the French text published by the District of France of
the Society of Saint Pius X; the wording is rendered, not quoted. The
gloss ``the exclusive celebration of the liturgy according to the books of
the \latin{editio typica} of 1962'' is Pietras's, art.\ cit., 179, not the
recourse's; the recourse itself speaks of the \latin{elemento specifico}
for which the society was founded and of ``the celebration of the ancient
and venerable form of the Roman liturgy.'' The distinction between the
document's word and its scholarly gloss is preserved here deliberately.}
That is not a quibble. In 2000 the Holy See refused to make
exclusivity a norm of the institute, and in 2003 it definitively approved
a constitutional text that does not assert it.

\emph{It incorporates \work{Ecclesia Dei} by reference.} The traditions to
be faithfully observed are those ``according to the dispositions of the
Motu proprio Ecclesia Dei.'' A constitution that defines its own
liturgical patrimony by reference to a papal act is exposed to whatever
happens to that act---a point that became acute in 2021, when
\work{Traditionis custodes} abrogated ``previous norms, instructions,
permissions, and customs'' not conforming to it, and again when the
Commission that \work{Ecclesia Dei} instituted was itself decommissioned.

\emph{It states an aim, not a permission.} The permission to use the 1962
books is in the erection decree and in the decree of 10 September 1988;
article 8 states what the Fraternity is for. The distinction is the
Fraternity's own best argument and its own greatest exposure: an aim
approved by the Apostolic See is part of the institute's patrimony under
canon 578 and enjoys the just autonomy of canon 586, but an aim cannot by
itself supply a faculty that the liturgical law has withdrawn.

\subsection{The rest of Part I}

The remaining published articles fill in what a reader needs to make
sense of later disputes. The Sacrifice of the Mass is ``at the heart of
the spirituality and the apostolate'' (art.~2). The Fraternity is under
the protection of the Blessed Virgin (art.~4) and the patronage of Saint
Peter, ``in order to express their gratitude, filial love, and loyalty to
the Supreme Pontiff'' (art.~5). It devotes itself to ``all the works of
priestly formation,'' with philosophical and theological studies
``founded on the principles and the method of Saint Thomas Aquinas''
(art.~10); formation follows the \latin{ratio studiorum} promulgated by
the Holy See, with a directory specifying the seminary curriculum in
conformity with the law (art.~11). Parochial ministry ``is a work to
which the Fraternity devotes itself if a bishop requests of it such
services,'' to be the object of agreements with diocesan bishops
``in order to permit the Fraternity to exercise its apostolate according
to its proper charisma'' (art.~13). And article 15 provides for
vocations of auxiliaries in conformity with canon 738 §\,2.

That last citation is worth a sentence. Canon 738 §\,2 is the canon that
subjects members of a society of apostolic life to the diocesan bishop in
public worship, the care of souls and the works of the apostolate. The
Fraternity's own fundamental code cites it, and article 13 makes
parochial ministry depend on a bishop's request and a written agreement.
Whatever ``pontifical right'' means in this institute's self-description,
it has never meant, in its own constitutions, a claim to work in a
diocese without its bishop.

\actrecord{Constitutions of the Priestly Fraternity of Saint Peter,
definitively approved 29 June 2003 (Prot.\ N.~108/2003); published excerpt
of Part I, articles 1--15.}{That the Fraternity's fundamental code
defines it as a clerical society of apostolic life of pontifical right;
professes fidelity to the Roman Pontiff with canon 590 §\,2; gives as its
object the sanctification of priests through the priesthood and the Mass;
gives as its particular aim the faithful observance of the liturgical and
disciplinary traditions according to the dispositions of \work{Ecclesia
Dei}; grounds its studies in Saint Thomas; and makes its parochial
ministry depend on a bishop's request and an agreement.}{Only fifteen
articles are public, in a translation published by the institute. The
articles on government, incorporation, formation and the chapter are not
public, and this account cannot check the reading of any of them,
including the article 19 invoked in the 2000 recourse.}
